\chapter{Experimental Framework}
\label{ch:benchmark}

Antes de presentar las mejoras propuestas sobre OVO, conviene describir el marco de trabajo
sobre el que se fundamentan: los datasets y las métricas con los que se evalúa el sistema, y
las herramientas creadas y utilizadas para medir sus resultados, observar su comportamiento y
provocar sus fallos de forma controlada y reproducible.

Todo este aparato responde a una exigencia de atribución: poder afirmar que una mejora se
debe a un cambio concreto y no al azar. Se presenta en tres piezas, cada una respondiendo a
una pregunta: con qué se juzga (Sección~\ref{sec:bench_setup}), cómo se provoca el fallo de
forma repetible (Sección~\ref{sec:bench_simslam}) y desde qué punto se parte
(Sección~\ref{sec:bench_baseline}).

\section{Dataset and metrics}
\label{sec:bench_setup}

La evaluación principal se realiza sobre Replica~\parencite{straub2019replica}, un conjunto
de escenas de interior reconstruidas con geometría densa y anotación semántica por punto.
Se usa el subconjunto habitual en la literatura de OVO: ocho escenas (office-0 a office-4 y
room-0 a room-2) y 51 clases. Su anotación densa por punto y sus
trayectorias limpias, sin el ruido de captura de un dataset real, permiten aislar el comportamiento
semántico del ruido geométrico.


% TODO: añadir cita de ScanNet a refs.bib (dai2017scannet) y ponerla aquí.
Junto a Replica se emplean también dos escenas de ScanNet, no como \emph{baseline} de
métricas reproducible sino como banco de prueba en condiciones más realistas: permiten
evaluar OVO con ORB-SLAM2 sobre tracking real, con el ruido de captura y la deriva que la
trayectoria limpia de Replica no tiene (Capítulos~\ref{ch:geometric_consistency}
y~\ref{ch:fusion}).

La calidad del mapa se mide en las dos familias de métricas que el
Capítulo~\ref{ch:preliminaries} anticipa, cada una atada a un eje distinto de fallo:

\begin{itemize}
  \item \textbf{Class-aware} (mIoU, mAcc). Comparan la etiqueta de cada punto con la del GT y
    miden la calidad \emph{semántica}. Se reportan por tertiles de frecuencia de clase
    (Head, Common, Tail), siguiendo la práctica de OpenNeRF: promediar sobre las 51 clases
    dejaría que las pocas clases muy frecuentes taparan el comportamiento en la cola larga,
    donde vive la mayor parte del vocabulario.
  \item \textbf{Class-agnostic} (AP de instancia, con sus variantes AP, AP$_{50}$, AP$_{25}$).
    Comparan la segmentación en instancias con la del GT \emph{sin mirar la clase}, y miden
    si el objeto se recuperó como una instancia correcta.
\end{itemize}

La parte semántica se resuelve con el segmentador SAM~2.1-l~\parencite{ravi2024sam2} y el
descriptor SigLIP ViT-SO400M~\parencite{zhai2023siglip}. Todo se ejecuta sobre una única GPU
NVIDIA RTX 4090 (24\,GB).

Para observar el sistema por dentro se ha construido, como parte de este trabajo, un visor
basado en Rerun. La visualización previa, en Open3D, mostraba la nube y permitía consultas por
lenguaje, pero no recorrer la construcción keyframe a keyframe ni rebobinarla. Rerun registra
el mapa sobre una línea temporal, con paneles 2D por keyframe (RGB, máscaras de SAM y
reproyecciones sobre la imagen), lo que permite avanzar y retroceder y rastrear una mala
segmentación hasta su \emph{origen temporal} (Figura~\ref{fig:rerun_tracking}). Aunque la evaluación
cuantitativa de los resultados recae en las métricas, es esta inspección visual, que presenta
de forma simultánea instancias, máscaras y trayectoria frame a frame, la que permite localizar
los fallos y su causa y, a partir de ahí, plantear mejoras.

\begin{figure}[H]
  \centering
  % PLACEHOLDER: sustituir por \includegraphics de una captura del visor Rerun (PNG, modo stream).
  \fbox{\parbox[c][3cm][c]{0.8\linewidth}{\centering\footnotesize\itshape
  Captura de Rerun: nube por instancia, paneles 2D por keyframe (RGB, máscaras de SAM,
  reproyecciones) y el origen temporal de una mala asignación.}}
  \caption{El visor Rerun, con el que se observa la construcción del mapa keyframe a keyframe.}
  \label{fig:rerun_tracking}
\end{figure}

\section{SimulatedSLAM: a backend with controlled perturbation}
\label{sec:bench_simslam}

Varias de las contribuciones de este trabajo estudian cómo afectan a la fusión de instancias
las imperfecciones de la geometría y la corrección brusca que introduce un loop closure.
Reproducirlas con un SLAM real es difícil, porque hay que esperar a que ocurran y no se
repiten igual entre ejecuciones. Por eso se añade un tercer backend, junto
al SLAM vanilla y a ORB-SLAM2~\parencite{murartal2017orbslam2}: SimulatedSLAM, que reconstruye la nube igual que
OVO-mapping a partir de la trayectoria \emph{ground truth}, pero con la capacidad de inyectar
saltos de pose de forma controlada.

Con la inyección apagada entrega las poses GT sin tocarlas y reproduce, por tanto, las cifras
de OVO-mapping. La Sección~\ref{sec:bench_baseline} lo confirma, y esa igualdad valida el
banco: si el sustrato limpio equivale a OVO-mapping, cualquier degradación posterior es
atribuible a la perturbación inyectada y no al backend de simulación.

Con la inyección encendida, SimulatedSLAM introduce en un keyframe elegido un salto brusco de
pose. El salto es una transformación rígida
$J = \left[\begin{smallmatrix} R & t \\ 0 & 1\end{smallmatrix}\right]$, con rotación $R$ y
traslación $t$, que se aplica a todas las poses a partir de ese instante: la pose estimada del
keyframe $n$ (la que el SLAM reporta y con la que se colocan sus puntos) deja de coincidir con
su pose real $T^{\mathrm{gt}}_n$ y pasa a ser $T^{\mathrm{est}}_n = J\,T^{\mathrm{gt}}_n$. Del
salto se eligen el keyframe de
disparo y sus magnitudes: cuánto traslada y cuántos grados gira. Lo que fija la semilla es la
orientación, la dirección de la traslación y el eje alrededor del cual gira, de modo que el
salto es aleatorio en orientación pero perfectamente reproducible. Según qué componente esté activa, el salto es
\emph{traslacional} ($R=I$), \emph{rotacional} ($t=0$) o \emph{mixto} (ambos). Al producirse el
salto, las instancias presentes en ese frame se reconstruyen desplazadas a otra zona del
espacio, generando duplicados por construcción. Con ellos se reproduce el error geométrico de
un SLAM real: al cerrarse el loop closure, esos duplicados se corrigen y quedan superpuestos
sobre su instancia original. Al final de la secuencia, SimulatedSLAM simula ese  loop closure devolviendo el mapa a su
\emph{ground truth}: fija cada keyframe a su pose real $T^{\mathrm{gt}}_n$ y arrastra
rígidamente sus puntos con la transformación de corrección $C_n$,
\begin{equation}
  \label{eq:oracle_correction}
  C_n = T^{\mathrm{gt}}_n\,(T^{\mathrm{est}}_n)^{-1}, \qquad p' = C_n\,p \quad \forall\, p \in \mathcal{P}_n,
\end{equation}
que lleva cada punto $p$ del keyframe (el conjunto $\mathcal{P}_n$, en coordenadas homogéneas) a
su posición corregida $p'$. Así se ajusta a la vez trayectoria y nube. En los keyframes
posteriores al salto $C_n = J^{-1}$: deshace exactamente el salto inyectado, y deja intactos los
anteriores. Es una corrección perfecta, instantánea y única, al contrario que la de un
SLAM real. El protocolo que la explota para puntuar la fusión se detalla en el
Capítulo~\ref{ch:fusion}.

\section{Baseline: reproducing OVO}
\label{sec:bench_baseline}

Para comprobar que los experimentos parten de un buen punto, se ejecuta OVO sobre el entorno
de pruebas con el objetivo de reproducir las cifras del artículo original y fijar el
\emph{baseline} de referencia. Se ejecuta también SimulatedSLAM, con la inyección apagada,
para verificar que reconstruye correctamente y reproduce ese mismo baseline. La
Tabla~\ref{tab:baseline_vs_paper} recoge ambas comprobaciones.

% Baseline reproduction vs. published OVO results (Replica, 51 classes).
% Semantic: frequency tertiles Head/Common/Tail (mIoU/mAcc, %).
% AP: class-agnostic instance AP (AP/AP50/AP25, %). Paper does not report it -> "--".
% Paper rows: OVO Table II. "(reprod.)" rows: our initial baseline, this work.
% ORB-SLAM2 (reprod.) is the mean of 2 runs (ORB-SLAM is non-deterministic;
%   observed spread ~2 mIoU / ~2 AP between runs).
% % Baseline reproduction vs. published OVO results (Replica, 51 classes).
% Semantic: frequency tertiles Head/Common/Tail (mIoU/mAcc, %).
% AP: class-agnostic instance AP (AP/AP50/AP25, %). Paper does not report it -> "--".
% Paper rows: OVO Table II. "(reprod.)" rows: our initial baseline, this work.
% ORB-SLAM2 (reprod.) is the mean of 2 runs (ORB-SLAM is non-deterministic;
%   observed spread ~2 mIoU / ~2 AP between runs).
% % Baseline reproduction vs. published OVO results (Replica, 51 classes).
% Semantic: frequency tertiles Head/Common/Tail (mIoU/mAcc, %).
% AP: class-agnostic instance AP (AP/AP50/AP25, %). Paper does not report it -> "--".
% Paper rows: OVO Table II. "(reprod.)" rows: our initial baseline, this work.
% ORB-SLAM2 (reprod.) is the mean of 2 runs (ORB-SLAM is non-deterministic;
%   observed spread ~2 mIoU / ~2 AP between runs).
% \input{figures/tab_baseline_vs_paper}
\begin{table}[H]
  \centering
  \setlength{\tabcolsep}{4pt}
  \renewcommand{\arraystretch}{1.15}
  \resizebox{\textwidth}{!}{%
  \begin{tabular}{l cc cc cc cc c ccc}
    \toprule
    & \multicolumn{2}{c}{All} & \multicolumn{2}{c}{Head}
    & \multicolumn{2}{c}{Common} & \multicolumn{2}{c}{Tail}
    & & \multicolumn{3}{c}{AP class-agnostic} \\
    \cmidrule(lr){2-3}\cmidrule(lr){4-5}\cmidrule(lr){6-7}\cmidrule(lr){8-9}\cmidrule(lr){11-13}
    Method & mIoU & mAcc & mIoU & mAcc & mIoU & mAcc & mIoU & mAcc & & AP & AP$_{50}$ & AP$_{25}$ \\
    \midrule
    OVO-mapping (paper)      & 27.0 & 39.1 & 45.0 & 59.9 & 25.1 & 38.5 & 11.0 & 18.8 & & -- & -- & -- \\
    \rowcolor{gray!12}
    OVO-mapping (reprod.)    & 27.6 & 40.1 & 44.5 & 59.2 & 26.4 & 42.2 & 12.0 & 19.1 & & 23.4 & 40.9 & 59.3 \\
    \rowcolor{gray!12}
    OVO-Simulated (reprod.)  & 27.1 & 40.2 & 44.3 & 59.0 & 25.7 & 42.1 & 11.2 & 19.5 & & 23.4 & 41.1 & 59.7 \\
    \midrule
    OVO-ORB-SLAM2 (paper)    & 25.6 & 39.0 & 43.0 & 59.1 & 21.6 & 38.3 & 12.1 & 19.6 & & -- & -- & -- \\
    \rowcolor{gray!12}
    OVO-ORB-SLAM2 (reprod.)  & 26.0 & 39.2 & 43.2 & 58.7 & 23.5 & 39.5 & 11.4 & 19.4 & & 22.9 & 40.1 & 60.4 \\
    \bottomrule
  \end{tabular}%
  }
  \caption{Published OVO results versus our reproduction of the \emph{baseline} on
    Replica (51 classes).}
  \label{tab:baseline_vs_paper}
\end{table}

\begin{table}[H]
  \centering
  \setlength{\tabcolsep}{4pt}
  \renewcommand{\arraystretch}{1.15}
  \resizebox{\textwidth}{!}{%
  \begin{tabular}{l cc cc cc cc c ccc}
    \toprule
    & \multicolumn{2}{c}{All} & \multicolumn{2}{c}{Head}
    & \multicolumn{2}{c}{Common} & \multicolumn{2}{c}{Tail}
    & & \multicolumn{3}{c}{AP class-agnostic} \\
    \cmidrule(lr){2-3}\cmidrule(lr){4-5}\cmidrule(lr){6-7}\cmidrule(lr){8-9}\cmidrule(lr){11-13}
    Method & mIoU & mAcc & mIoU & mAcc & mIoU & mAcc & mIoU & mAcc & & AP & AP$_{50}$ & AP$_{25}$ \\
    \midrule
    OVO-mapping (paper)      & 27.0 & 39.1 & 45.0 & 59.9 & 25.1 & 38.5 & 11.0 & 18.8 & & -- & -- & -- \\
    \rowcolor{gray!12}
    OVO-mapping (reprod.)    & 27.6 & 40.1 & 44.5 & 59.2 & 26.4 & 42.2 & 12.0 & 19.1 & & 23.4 & 40.9 & 59.3 \\
    \rowcolor{gray!12}
    OVO-Simulated (reprod.)  & 27.1 & 40.2 & 44.3 & 59.0 & 25.7 & 42.1 & 11.2 & 19.5 & & 23.4 & 41.1 & 59.7 \\
    \midrule
    OVO-ORB-SLAM2 (paper)    & 25.6 & 39.0 & 43.0 & 59.1 & 21.6 & 38.3 & 12.1 & 19.6 & & -- & -- & -- \\
    \rowcolor{gray!12}
    OVO-ORB-SLAM2 (reprod.)  & 26.0 & 39.2 & 43.2 & 58.7 & 23.5 & 39.5 & 11.4 & 19.4 & & 22.9 & 40.1 & 60.4 \\
    \bottomrule
  \end{tabular}%
  }
  \caption{Published OVO results versus our reproduction of the \emph{baseline} on
    Replica (51 classes).}
  \label{tab:baseline_vs_paper}
\end{table}

\begin{table}[H]
  \centering
  \setlength{\tabcolsep}{4pt}
  \renewcommand{\arraystretch}{1.15}
  \resizebox{\textwidth}{!}{%
  \begin{tabular}{l cc cc cc cc c ccc}
    \toprule
    & \multicolumn{2}{c}{All} & \multicolumn{2}{c}{Head}
    & \multicolumn{2}{c}{Common} & \multicolumn{2}{c}{Tail}
    & & \multicolumn{3}{c}{AP class-agnostic} \\
    \cmidrule(lr){2-3}\cmidrule(lr){4-5}\cmidrule(lr){6-7}\cmidrule(lr){8-9}\cmidrule(lr){11-13}
    Method & mIoU & mAcc & mIoU & mAcc & mIoU & mAcc & mIoU & mAcc & & AP & AP$_{50}$ & AP$_{25}$ \\
    \midrule
    OVO-mapping (paper)      & 27.0 & 39.1 & 45.0 & 59.9 & 25.1 & 38.5 & 11.0 & 18.8 & & -- & -- & -- \\
    \rowcolor{gray!12}
    OVO-mapping (reprod.)    & 27.6 & 40.1 & 44.5 & 59.2 & 26.4 & 42.2 & 12.0 & 19.1 & & 23.4 & 40.9 & 59.3 \\
    \rowcolor{gray!12}
    OVO-Simulated (reprod.)  & 27.1 & 40.2 & 44.3 & 59.0 & 25.7 & 42.1 & 11.2 & 19.5 & & 23.4 & 41.1 & 59.7 \\
    \midrule
    OVO-ORB-SLAM2 (paper)    & 25.6 & 39.0 & 43.0 & 59.1 & 21.6 & 38.3 & 12.1 & 19.6 & & -- & -- & -- \\
    \rowcolor{gray!12}
    OVO-ORB-SLAM2 (reprod.)  & 26.0 & 39.2 & 43.2 & 58.7 & 23.5 & 39.5 & 11.4 & 19.4 & & 22.9 & 40.1 & 60.4 \\
    \bottomrule
  \end{tabular}%
  }
  \caption{Published OVO results versus our reproduction of the \emph{baseline} on
    Replica (51 classes).}
  \label{tab:baseline_vs_paper}
\end{table}


La reproducción recupera las cifras del artículo dentro de un par de puntos, con la clase
Common algo por encima (el artículo no reporta AP de instancia, de ahí los guiones); la fila
OVO-ORB-SLAM2 es la media de cuatro ejecuciones, dado que ORB-SLAM es no determinista, con una
dispersión de unos dos puntos en mIoU y AP. La fila OVO-Simulated, la relevante para este
trabajo, coincide casi exactamente con OVO-mapping (All $27.1$ frente a $27.6$ mIoU); la
pequeña diferencia proviene de que SimulatedSLAM selecciona sus keyframes por distancia y
rotación, y reconstruye por tanto a partir de menos frames que OVO-mapping, no de un error
introducido por la simulación. Esa coincidencia lo valida como el \emph{baseline} de ruido
cero, la fila contra la que se mide todo lo que sigue.
